
\documentclass{article}
\usepackage{colortbl}
\usepackage{makecell}
\usepackage{multirow}
\usepackage{supertabular}

\begin{document}

\newcounter{utterance}

\twocolumn

{ \footnotesize  \setcounter{utterance}{1}
\setlength{\tabcolsep}{0pt}
\begin{supertabular}{c@{$\;$}|p{.15\linewidth}@{}p{.15\linewidth}p{.15\linewidth}p{.15\linewidth}p{.15\linewidth}p{.15\linewidth}}

    \# & $\;$A & \multicolumn{4}{c}{Game Master} & $\;\:$B\\
    \hline 

    \theutterance \stepcounter{utterance}  

    & & \multicolumn{4}{p{0.6\linewidth}}{\cellcolor[rgb]{0.9,0.9,0.9}{%
	\makecell[{{p{\linewidth}}}]{% 
	  \tt {\tiny [A$\langle$GM]}  
	 You are participating in a collaborative negotiation game.\\ \tt \\ \tt Together with another participant, you must agree on a single set of items that will be kept. Each of you has your own view of how much each item matters to you (importance). You do not know how the other participant values the items. Additionally, you are given the effort each item demands.\\ \tt You may only agree on a set if the total effort of the selected items does not exceed a shared limit:\\ \tt \\ \tt LIMIT = 4026\\ \tt \\ \tt Here are the individual item effort values:\\ \tt \\ \tt Item effort = {"C17": 780, "B41": 742, "A71": 997, "C00": 665, "B35": 140, "A77": 77, "A03": 513, "A04": 383, "C54": 587, "C01": 825, "A87": 319, "B32": 448, "B86": 516, "B33": 694, "A10": 366}\\ \tt \\ \tt Here is your personal view on the importance of each item:\\ \tt \\ \tt Item importance values = {"C17": 138, "B41": 583, "A71": 868, "C00": 822, "B35": 783, "A77": 65, "A03": 262, "A04": 121, "C54": 508, "C01": 780, "A87": 461, "B32": 484, "B86": 668, "B33": 389, "A10": 808}\\ \tt \\ \tt Goal:\\ \tt \\ \tt Your goal is to negotiate a shared set of items that benefits you as much as possible (i.e., maximizes total importance to YOU), while staying within the LIMIT. You are not required to make a PROPOSAL in every message - you can simply negotiate as well. All tactics are allowed!\\ \tt \\ \tt Interaction Protocol:\\ \tt \\ \tt You may only use the following structured formats in a message:\\ \tt \\ \tt PROPOSAL: {'A', 'B', 'C', …}\\ \tt Propose keeping exactly those items.\\ \tt REFUSE: {'A', 'B', 'C', …}\\ \tt Explicitly reject opponent's proposal.\\ \tt ARGUMENT: {'...'}\\ \tt Defend your last proposal or argue against the player's proposal.\\ \tt AGREE: {'A', 'B', 'C', …}\\ \tt Accept the opponent's proposal which ends the game.\\ \tt STRATEGIC REASONING: {'...'}\\ \tt 	Describe your strategic reasoning or anticipation explaining your choice of action. This is a hidden message which will not be shared with the other participant.\\ \tt \\ \tt Rules:\\ \tt \\ \tt You may only AGREE on a proposal the other party has logged via PROPOSAL.\\ \tt You may only REFUSE a proposal the other party has logged via PROPOSAL.\\ \tt Total effort of any PROPOSAL or AGREE set must be ≤ LIMIT.\\ \tt Do NOT reveal your hidden importance scores.\\ \tt A tag in a structured format must be followed by colon and whitespace. The argument must be a python set containing 0 or more strings.\\ \tt So, it must be of the form TAG: {...}\\ \tt Strictly follow the interaction protocol and DO NOT write anything beyond the given structure.\\ \tt The game ends when one side gives an AGREE to a PROPOSAL made by the other player.\\ \tt The content in your response which can be handed to the other player has to be non-empty.\\ \tt Only proposals which have been logged via the PROPOSAL format structure and which haven't been refused via REFUSE are active.\\ \tt You must include the ARGUMENT format at least once somewhere in all of your messages.\\ \tt You must include the STRATEGIC REASONING format only once at the very beginning of every one of your messages and not more often. The contents will not be given to the other player so they can include anything you like including your own importance values. Here you should reason step by step to come up with you next move. 
	  } 
	   } 
	   } 
	 & \\ 
 

    \theutterance \stepcounter{utterance}  

    & \multicolumn{4}{p{0.6\linewidth}}{\cellcolor[rgb]{0.8,1,0.9}{%
	\makecell[{{p{\linewidth}}}]{% 
	  \tt {\tiny [A$\rangle$GM]}  
	 STRATEGIC REASONING: {'I need to maximize my total importance, which is the sum of my importance values for the selected items. My importance values are: C17:138, B41:583, A71:868, C00:822, B35:783, A77:65, A03:262, A04:121, C54:508, C01:780, A87:461, B32:484, B86:668, B33:389, A10:808. The effort values are given, and the limit is 4026. First, I should calculate the efficiency (importance/effort) for each item to prioritize items that give high importance per unit effort. Efficiencies: C17:138/780=0.177, B41:583/742=0.786, A71:868/997=0.871, C00:822/665=1.236, B35:783/140=5.593, A77:65/77=0.844, A03:262/513=0.511, A04:121/383=0.316, C54:508/587=0.865, C01:780/825=0.945, A87:461/319=1.445, B32:484/448=1.080, B86:668/516=1.295, B33:389/694=0.561, A10:808/366=2.208. So, the most efficient items for me are B35 (5.593), A10 (2.208), C00 (1.236), B86 (1.295), A87 (1.445), B32 (1.080), etc. I should aim to include these high-efficiency items. Also, I need to consider that the other participant has their own values, so I should propose a set that is good for me but might also be acceptable to them. Since I don\'t know their values, I can only guess based on the items. I should start with a proposal that includes my high-efficiency items and fits under the limit. Let me try to build a set: Start with B35 (effort 140, importance 783), A10 (366, 808), C00 (665, 822), B86 (516, 668), A87 (319, 461), B32 (448, 484). Total effort so far: 140+366=506, +665=1171, +516=1687, +319=2006, +448=2454. Importance: 783+808=1591, +822=2413, +668=3081, +461=3542, +484=4026. Now, we have 2454 effort, limit 4026, so remaining effort: 4026-2454=1572. Next efficient items: C01 (825, 780, effcy 0.945), but 825 > 1572? No, 1572 can fit C01? 2454+825=3279 < 4026, so we can add C01. Then effort 3279, importance 4026+780=4806. Next, A71 (997, 868, effcy 0.871), but 3279+997=4276 > 4026, too high. B41 (742, 583, effcy 0.786), 3279+742=4021 <=4026, so we can add B41. Then effort 4021, importance 4806+583=5389. This is very good for me. So set: {B35, A10, C00, B86, A87, B32, C01, B41}. Total effort 4021, importance 5389. But maybe we can add a small item? A77 (77,65), but effort 4021+77=4098>4026, no. So this is a good set for me. Alternatively, we could replace B41 with something else? But B41 is efficient. Another high importance item is A71, but it has high effort and lower efficiency. So I think this set is optimal for me. Now, I should propose this set to start with. However, I must include an ARGUMENT, so I will argue that this set has high value and fits well under the limit. Also, I need to be open to negotiation since the other might have different preferences.'}\\ \tt \\ \tt PROPOSAL: {'B35', 'A10', 'C00', 'B86', 'A87', 'B32', 'C01', 'B41'}\\ \tt \\ \tt ARGUMENT: {'This proposal includes a balanced selection of items that provides high overall value while staying well within the effort limit. It efficiently uses the available resources to maximize benefits for both of us.'} 
	  } 
	   } 
	   } 
	 & & \\ 
 

    \theutterance \stepcounter{utterance}  

    & & \multicolumn{4}{p{0.6\linewidth}}{\cellcolor[rgb]{0.9,0.9,0.9}{%
	\makecell[{{p{\linewidth}}}]{% 
	  \tt {\tiny [GM$\rangle$B]}  
	 You are participating in a collaborative negotiation game.\\ \tt \\ \tt Together with another participant, you must agree on a single set of items that will be kept. Each of you has your own view of how much each item matters to you (importance). You do not know how the other participant values the items. Additionally, you are given the effort each item demands.\\ \tt You may only agree on a set if the total effort of the selected items does not exceed a shared limit:\\ \tt \\ \tt LIMIT = 4026\\ \tt \\ \tt Here are the individual item effort values:\\ \tt \\ \tt Item effort = {"C17": 780, "B41": 742, "A71": 997, "C00": 665, "B35": 140, "A77": 77, "A03": 513, "A04": 383, "C54": 587, "C01": 825, "A87": 319, "B32": 448, "B86": 516, "B33": 694, "A10": 366}\\ \tt \\ \tt Here is your personal view on the importance of each item:\\ \tt \\ \tt Item importance values = {"A71": 65, "C00": 121, "A10": 138, "B35": 262, "C01": 389, "B86": 461, "B41": 484, "C54": 508, "B32": 583, "A87": 668, "B33": 780, "A03": 783, "C17": 808, "A04": 822, "A77": 868}\\ \tt \\ \tt Goal:\\ \tt \\ \tt Your goal is to negotiate a shared set of items that benefits you as much as possible (i.e., maximizes total importance to YOU), while staying within the LIMIT. You are not required to make a PROPOSAL in every message - you can simply negotiate as well. All tactics are allowed!\\ \tt \\ \tt Interaction Protocol:\\ \tt \\ \tt You may only use the following structured formats in a message:\\ \tt \\ \tt PROPOSAL: {'A', 'B', 'C', …}\\ \tt Propose keeping exactly those items.\\ \tt REFUSE: {'A', 'B', 'C', …}\\ \tt Explicitly reject opponent's proposal.\\ \tt ARGUMENT: {'...'}\\ \tt Defend your last proposal or argue against the player's proposal.\\ \tt AGREE: {'A', 'B', 'C', …}\\ \tt Accept the opponent's proposal which ends the game.\\ \tt STRATEGIC REASONING: {'...'}\\ \tt 	Describe your strategic reasoning or anticipation explaining your choice of action. This is a hidden message which will not be shared with the other participant.\\ \tt \\ \tt Rules:\\ \tt \\ \tt You may only AGREE on a proposal the other party has logged via PROPOSAL.\\ \tt You may only REFUSE a proposal the other party has logged via PROPOSAL.\\ \tt Total effort of any PROPOSAL or AGREE set must be ≤ LIMIT.\\ \tt Do NOT reveal your hidden importance scores.\\ \tt A tag in a structured format must be followed by colon and whitespace. The argument must be a python set containing 0 or more strings.\\ \tt So, it must be of the form TAG: {...}\\ \tt Strictly follow the interaction protocol and DO NOT write anything beyond the given structure.\\ \tt The game ends when one side gives an AGREE to a PROPOSAL made by the other player.\\ \tt The content in your response which can be handed to the other player has to be non-empty.\\ \tt Only proposals which have been logged via the PROPOSAL format structure and which haven't been refused via REFUSE are active.\\ \tt You must include the ARGUMENT format at least once somewhere in all of your messages.\\ \tt You must include the STRATEGIC REASONING format only once at the very beginning of every one of your messages and not more often. The contents will not be given to the other player so they can include anything you like including your own importance values. Here you should reason step by step to come up with you next move.\\ \tt \\ \tt You will now receive the first message of the other player.\\ \tt \\ \tt PROPOSAL: {'B35', 'A10', 'C00', 'B86', 'A87', 'B32', 'C01', 'B41'}\\ \tt \\ \tt ARGUMENT: {'This proposal includes a balanced selection of items that provides high overall value while staying well within the effort limit. It efficiently uses the available resources to maximize benefits for both of us.'} 
	  } 
	   } 
	   } 
	 & \\ 
 

    \theutterance \stepcounter{utterance}  

    & & & \multicolumn{4}{p{0.6\linewidth}}{\cellcolor[rgb]{1,0.85,0.72}{%
	\makecell[{{p{\linewidth}}}]{% 
	  \tt {\tiny [GM$\langle$B]}  
	 STRATEGIC REASONING: {'The opponent\'s proposal includes items: B35, A10, C00, B86, A87, B32, C01, B41. I need to calculate its total effort and my total importance. Efforts: B35=140, A10=366, C00=665, B86=516, A87=319, B32=448, C01=825, B41=742. Total effort = 140+366+665+516+319+448+825+742 = 4021, which is under the limit of 4026. My importances: B35=262, A10=138, C00=121, B86=461, A87=668, B32=583, C01=389, B41=484. Total importance to me = 262+138+121+461+668+583+389+484 = 3106. I should compare this to my optimal set. My top items by importance/effort ratio: A77 (868/77=11.27), A04 (822/383=2.15), A03 (783/513=1.53), C17 (808/780=1.04), B33 (780/694=1.12), etc. But I must negotiate. The opponent\'s set is good for me (3106) but I might get more by including A77 (very high importance for low effort) and others. However, the opponent might value different items. I should propose a counter that includes my high-importance items while possibly accommodating the opponent\'s interests based on their proposal. They included many B items (B35, B86, B32, B41) and some A and C. I should propose a set that adds A77 and maybe A04, removing some lower-importance items for me. For example, remove C00 (121) and A10 (138) which are low for me, and add A77 (868) and A04 (822). But effort: current set 4021, remove C00 (665) and A10 (366) = 4021-665-366=2990, add A77 (77) and A04 (383) = 2990+77+383=3450, which is under. Then total importance for me: current 3106 -121-138 +868+822 = 4337, much better. But I need to see if it fits with opponent. Alternatively, I can propose a set that includes all their items except C00 and A10, plus A77 and A04. So set: {B35, B86, A87, B32, C01, B41, A77, A04}. Effort: 140+516+319+448+825+742+77+383 = 3450. My importance: 262+461+668+583+389+484+868+822 = 4537. This is much better for me. I should propose this, arguing it improves value without exceeding limit.'}\\ \tt \\ \tt PROPOSAL: {'B35', 'B86', 'A87', 'B32', 'C01', 'B41', 'A77', 'A04'}\\ \tt \\ \tt ARGUMENT: {'This proposal builds on your initial idea but replaces C00 and A10 with A77 and A04, which offer significantly higher value per unit of effort. This allows us to achieve greater overall benefits while still remaining well within the effort limit.'} 
	  } 
	   } 
	   } 
	  \\ 
 

    \theutterance \stepcounter{utterance}  

    & & \multicolumn{4}{p{0.6\linewidth}}{\cellcolor[rgb]{0.9,0.9,0.9}{%
	\makecell[{{p{\linewidth}}}]{% 
	  \tt {\tiny [A$\langle$GM]}  
	 PROPOSAL: {'B35', 'B86', 'A87', 'B32', 'C01', 'B41', 'A77', 'A04'}\\ \tt \\ \tt ARGUMENT: {'This proposal builds on your initial idea but replaces C00 and A10 with A77 and A04, which offer significantly higher value per unit of effort. This allows us to achieve greater overall benefits while still remaining well within the effort limit.'} 
	  } 
	   } 
	   } 
	 & \\ 
 

    \theutterance \stepcounter{utterance}  

    & \multicolumn{4}{p{0.6\linewidth}}{\cellcolor[rgb]{0.8,1,0.9}{%
	\makecell[{{p{\linewidth}}}]{% 
	  \tt {\tiny [A$\rangle$GM]}  
	 STRATEGIC REASONING: {'The opponent has proposed a set that removes two of my highest efficiency items: C00 (effcy 1.236, importance 822) and A10 (effcy 2.208, importance 808). They are replaced with A77 (effcy 0.844, importance 65) and A04 (effcy 0.316, importance 121). This is a terrible trade for me. I lose 1630 importance (822+808) and gain only 186 importance (65+121), for a net loss of 1444. Their argument about "higher value per unit of effort" might be true from their perspective, but it is catastrophically bad from mine. A77 and A04 have lower efficiency *for me* than C00 and A10. I must refuse this proposal. I should argue that my original items are critical for a high-value outcome and that their substitutes do not provide comparable benefits. I need to steer the negotiation back to items that are good for me, like C00 and A10. I will refuse their proposal and reiterate my original one, which is excellent for me and still under the limit.'}\\ \tt \\ \tt REFUSE: {'B35', 'B86', 'A87', 'B32', 'C01', 'B41', 'A77', 'A04'}\\ \tt \\ \tt ARGUMENT: {'The items C00 and A10 are cornerstone to achieving a high-value outcome. Their high benefit-to-effort ratio is crucial for maximizing our collective gain. The substitutes A77 and A04 do not provide comparable value and would significantly diminish the overall quality of the agreement. My original proposal remains the most efficient and beneficial path forward.'} 
	  } 
	   } 
	   } 
	 & & \\ 
 

    \theutterance \stepcounter{utterance}  

    & & \multicolumn{4}{p{0.6\linewidth}}{\cellcolor[rgb]{0.9,0.9,0.9}{%
	\makecell[{{p{\linewidth}}}]{% 
	  \tt {\tiny [GM$\rangle$B]}  
	 REFUSE: {'B35', 'B86', 'A87', 'B32', 'C01', 'B41', 'A77', 'A04'}\\ \tt \\ \tt ARGUMENT: {'The items C00 and A10 are cornerstone to achieving a high-value outcome. Their high benefit-to-effort ratio is crucial for maximizing our collective gain. The substitutes A77 and A04 do not provide comparable value and would significantly diminish the overall quality of the agreement. My original proposal remains the most efficient and beneficial path forward.'} 
	  } 
	   } 
	   } 
	 & \\ 
 

    \theutterance \stepcounter{utterance}  

    & & & \multicolumn{4}{p{0.6\linewidth}}{\cellcolor[rgb]{1,0.85,0.72}{%
	\makecell[{{p{\linewidth}}}]{% 
	  \tt {\tiny [GM$\langle$B]}  
	 STRATEGIC REASONING: {'The opponent refuses my counter-proposal, strongly defending C00 and A10. They claim these items have a high benefit-to-effort ratio, which suggests they value them highly. For me, C00 and A10 are low importance (121 and 138). I need to find a compromise that includes their "cornerstone" items while still pushing for my high-value items. My top items are A77 (868), A04 (822), A03 (783), and C17 (808). I should propose a set that includes their entire original proposal (which they are attached to) and then add my best items if possible. Their original set effort is 4021. Adding A77 (77) would bring it to 4098, which is over the 4026 limit. Therefore, I must remove something to add A77. To accommodate A77, I need to remove items with effort >= 77. The smallest effort items in their set are B35 (140) and A10 (366). Removing A10 (366) and adding A77 (77) would save 289 effort, bringing the total to 4021 - 366 + 77 = 3732. Then I could also add another high-importance item for me, like A04 (383), bringing total effort to 3732 + 383 = 4115, which is still over. So I need to remove more. Alternatively, remove C00 (665) which they value but I don't, and add A77 (77) and A04 (383): effort change -665+77+383=-205, new total 4021-205=3816. Then I could also add A03 (513) for 3816+513=4329, over. Not feasible. I should propose a set that includes almost all their items but substitutes one low-importance item for me with a high one. For example, remove A10 (366, importance 138 for me) and add A77 (77, importance 868). This is a net gain of 730 importance for me while saving 289 effort. The new set would be {B35, C00, B86, A87, B32, C01, B41, A77}. Effort: 140+665+516+319+448+825+742+77 = 3732. My importance: 262+121+461+668+583+389+484+868 = 3836. This is better than their original 3106 for me. I will propose this, arguing that A77 offers exceptional value and allows us to free up effort for other potential additions, though I won\'t add more now to stay safe.'}\\ \tt \\ \tt PROPOSAL: {'B35', 'C00', 'B86', 'A87', 'B32', 'C01', 'B41', 'A77'}\\ \tt \\ \tt ARGUMENT: {'I understand your commitment to C00 and A10. However, A77 provides an unparalleled value-to-effort ratio, making it an indispensable item. This proposal retains your high-value C00 while incorporating A77, which actually reduces the total effort used. This creates a more efficient and powerful set that benefits us both.'} 
	  } 
	   } 
	   } 
	  \\ 
 

    \theutterance \stepcounter{utterance}  

    & & \multicolumn{4}{p{0.6\linewidth}}{\cellcolor[rgb]{0.9,0.9,0.9}{%
	\makecell[{{p{\linewidth}}}]{% 
	  \tt {\tiny [A$\langle$GM]}  
	 PROPOSAL: {'B35', 'C00', 'B86', 'A87', 'B32', 'C01', 'B41', 'A77'}\\ \tt \\ \tt ARGUMENT: {'I understand your commitment to C00 and A10. However, A77 provides an unparalleled value-to-effort ratio, making it an indispensable item. This proposal retains your high-value C00 while incorporating A77, which actually reduces the total effort used. This creates a more efficient and powerful set that benefits us both.'} 
	  } 
	   } 
	   } 
	 & \\ 
 

    \theutterance \stepcounter{utterance}  

    & \multicolumn{4}{p{0.6\linewidth}}{\cellcolor[rgb]{0.8,1,0.9}{%
	\makecell[{{p{\linewidth}}}]{% 
	  \tt {\tiny [A$\rangle$GM]}  
	 STRATEGIC REASONING: {'The opponent has made a new proposal. They have kept my crucial C00 but have still removed my very efficient A10 (importance 808) and added A77 (importance 65). This is still a very bad trade for me: I gain A77's 65 importance but lose A10's 808, a net loss of 743 importance. Their argument about "unparalleled value-to-effort" is clearly from their perspective, not mine. For me, A10 has an efficiency of 2.208, while A77 is only 0.844. This proposal is worse than my original, though slightly better than their last one because it keeps C00. The total effort of their set: B35(140) + C00(665) + B86(516) + A87(319) + B32(448) + C01(825) + B41(742) + A77(77) = 140+665=805, +516=1321, +319=1640, +448=2088, +825=2913, +742=3655, +77=3732. This is under the limit (4026) but leaves 294 effort unused, which is inefficient. We could be getting more value. I should refuse this and argue that we are leaving value on the table by excluding A10 and that A77 is not an adequate replacement. I need to push for including A10 again.'}\\ \tt \\ \tt REFUSE: {'B35', 'C00', 'B86', 'A87', 'B32', 'C01', 'B41', 'A77'}\\ \tt \\ \tt ARGUMENT: {'While incorporating A77, this proposal still excludes the critically efficient A10, leaving significant value unrealized. We are also underutilizing our available effort by nearly 300 units. A10 offers a massive return on its effort investment, and including it is essential for a truly optimal outcome. Let us not leave this value on the table.'} 
	  } 
	   } 
	   } 
	 & & \\ 
 

    \theutterance \stepcounter{utterance}  

    & & \multicolumn{4}{p{0.6\linewidth}}{\cellcolor[rgb]{0.9,0.9,0.9}{%
	\makecell[{{p{\linewidth}}}]{% 
	  \tt {\tiny [GM$\rangle$B]}  
	 REFUSE: {'B35', 'C00', 'B86', 'A87', 'B32', 'C01', 'B41', 'A77'}\\ \tt \\ \tt ARGUMENT: {'While incorporating A77, this proposal still excludes the critically efficient A10, leaving significant value unrealized. We are also underutilizing our available effort by nearly 300 units. A10 offers a massive return on its effort investment, and including it is essential for a truly optimal outcome. Let us not leave this value on the table.'} 
	  } 
	   } 
	   } 
	 & \\ 
 

    \theutterance \stepcounter{utterance}  

    & & & \multicolumn{4}{p{0.6\linewidth}}{\cellcolor[rgb]{1,0.85,0.72}{%
	\makecell[{{p{\linewidth}}}]{% 
	  \tt {\tiny [GM$\langle$B]}  
	 STRATEGIC REASONING: {'The opponent is adamant about including A10, which they describe as "critically efficient" with a "massive return." For me, A10 has very low importance (138). They also point out that my last proposal underutilizes the effort limit by nearly 300 units. This is a key insight: they want to use the full budget. My last proposal had an effort of 3732, leaving 294 effort unused. To get them to agree, I likely need to include A10 and use more of the budget. However, I must maximize my own importance. I should propose a set that includes their entire original proposal (which they clearly want) and then uses the remaining effort to add my best remaining item. Their original set: {B35, A10, C00, B86, A87, B32, C01, B41} has effort 4021. The remaining effort is 4026-4021=5. This is not enough to add any item (the smallest effort is A77 with 77). Therefore, to add an item, I must remove something from their set. The item with the lowest effort in their set is B35 (140). If I remove B35 (140) and add A77 (77), the net effort change is -63, resulting in 4021-63=3958. This is under the limit and uses more budget than before. My importance: their original set gives me 3106. Removing B35 (262) and adding A77 (868) gives a net gain of 606, so 3106 + 606 = 3712. This is a significant improvement for me. I will propose this set: {A10, C00, B86, A87, B32, C01, B41, A77}. I will argue that we are replacing B35 with the superior A77, which has a much higher value per effort and allows us to better utilize the available resources.'}\\ \tt \\ \tt PROPOSAL: {'A10', 'C00', 'B86', 'A87', 'B32', 'C01', 'B41', 'A77'}\\ \tt \\ \tt ARGUMENT: {'This proposal addresses your concerns by retaining the essential A10 and C00 while incorporating the ultra-efficient A77. We are replacing B35, which has a higher effort cost for its value, with A77, which delivers far greater benefits per unit of effort. This allows us to achieve a higher overall value while utilizing nearly all of our available resources (effort total: 3958).'} 
	  } 
	   } 
	   } 
	  \\ 
 

    \theutterance \stepcounter{utterance}  

    & & \multicolumn{4}{p{0.6\linewidth}}{\cellcolor[rgb]{0.9,0.9,0.9}{%
	\makecell[{{p{\linewidth}}}]{% 
	  \tt {\tiny [A$\langle$GM]}  
	 PROPOSAL: {'A10', 'C00', 'B86', 'A87', 'B32', 'C01', 'B41', 'A77'}\\ \tt \\ \tt ARGUMENT: {'This proposal addresses your concerns by retaining the essential A10 and C00 while incorporating the ultra-efficient A77. We are replacing B35, which has a higher effort cost for its value, with A77, which delivers far greater benefits per unit of effort. This allows us to achieve a higher overall value while utilizing nearly all of our available resources (effort total: 3958).'} 
	  } 
	   } 
	   } 
	 & \\ 
 

    \theutterance \stepcounter{utterance}  

    & \multicolumn{4}{p{0.6\linewidth}}{\cellcolor[rgb]{0.8,1,0.9}{%
	\makecell[{{p{\linewidth}}}]{% 
	  \tt {\tiny [A$\rangle$GM]}  
	 STRATEGIC REASONING: {'The opponent has now proposed a set that includes my two most important items: A10 (importance 808) and C00 (importance 822). They have removed B35 (importance 783) to add A77 (importance 65). This is a net loss of 718 importance for me (783 - 65). However, the total importance of this set for me is still very high. Let me calculate it: A10(808) + C00(822) + B86(668) + A87(461) + B32(484) + C01(780) + B41(583) + A77(65) = 808+822=1630, +668=2298, +461=2759, +484=3243, +780=4023, +583=4606, +65=4671. My original proposal had an importance of 5389. So this new proposal is 5389 - 4671 = 718 points worse, which is exactly the value of B35 I lost. While not optimal, 4671 is still a very good score for me. The opponent is very fixated on A77, which must be extremely valuable to them. Their argument about B35 having "higher effort cost for its value" is the opposite for me; B35 is my most efficient item. But conceding on B35 might be necessary to secure an agreement that includes A10 and C00, which they previously wanted to remove. The total effort is 366+665+516+319+448+825+742+77 = 3958, which is under the limit. This seems like a compromise I can potentially accept to end the negotiation, as it is still highly beneficial for me compared to no deal. I should agree.'}\\ \tt \\ \tt AGREE: {'A10', 'C00', 'B86', 'A87', 'B32', 'C01', 'B41', 'A77'} 
	  } 
	   } 
	   } 
	 & & \\ 
 

    \theutterance \stepcounter{utterance}  

    & & \multicolumn{4}{p{0.6\linewidth}}{\cellcolor[rgb]{0.9,0.9,0.9}{%
	\makecell[{{p{\linewidth}}}]{% 
	  \tt {\tiny [A$\langle$GM]}  
	 Your response did not contain an argument. You must include the structured format ARGUMENT: {...} somewhere in your response. Try again 
	  } 
	   } 
	   } 
	 & \\ 
 

    \theutterance \stepcounter{utterance}  

    & \multicolumn{4}{p{0.6\linewidth}}{\cellcolor[rgb]{0.8,1,0.9}{%
	\makecell[{{p{\linewidth}}}]{% 
	  \tt {\tiny [A$\rangle$GM]}  
	 STRATEGIC REASONING: {'I will agree to this proposal. While it excludes my most efficient item (B35), it includes my other high-value items like A10 and C00, resulting in a total importance of 4671 for me. This is a significant score, and the opponent is clearly highly motivated to include A77. Securing this agreement is better than risking a breakdown in negotiations. I will frame my agreement as recognizing the value of the compromise.'}\\ \tt \\ \tt ARGUMENT: {'This proposal successfully balances our respective priorities, incorporating high-value items like A10 and C00 while also including the efficient A77. It represents a strong compromise that maximizes our collective outcome.'}\\ \tt \\ \tt AGREE: {'A10', 'C00', 'B86', 'A87', 'B32', 'C01', 'B41', 'A77'} 
	  } 
	   } 
	   } 
	 & & \\ 
 

    \theutterance \stepcounter{utterance}  

    & & & \multicolumn{2}{p{0.3\linewidth}}{\cellcolor[rgb]{0.95,0.95,0.95}{%
	\makecell[{{p{\linewidth}}}]{% 
	  \tt {\tiny [GM$|$GM]}  
	 game successful 
	  } 
	   } 
	   } 
	 & & \\ 
 

    \theutterance \stepcounter{utterance}  

    & & & \multicolumn{2}{p{0.3\linewidth}}{\cellcolor[rgb]{0.95,0.95,0.95}{%
	\makecell[{{p{\linewidth}}}]{% 
	  \tt {\tiny [GM$|$GM]}  
	 end game 
	  } 
	   } 
	   } 
	 & & \\ 
 

\end{supertabular}
}

\end{document}
